\section{Physical ratio gates and resolved Wilson laws}
\label{sec:ratio}

This section uses three precise inputs from Paper 51~\cite{p51}: its finite-rank
tree transfer formula, its compact-product joint symbol theorem, and
its signed whole-alcove prefix bound. The extension below removes the
matched-product gates from the transition regions, while retaining
explicit negative-center profiles and literal color-ratio marks.
It does not assert the fully unmarked next logarithmic coefficient.

\subsection{The resolved angular measure}
Let $G$ be a fixed tree with its unique perfect matching
$M=\{a_jb_j\}_{j=1}^d$, so $N=2d$. An unmatched edge $a_i b_j$
defines an arrow $i\to j$. Choose a subset $E'$ of these arrows,
put $p=d-|E'|$, and contract $E'$ in the matching tree. The remaining
arrows define the order polytope $\Omega_{E'}\subset[0,1]^p$.
Write $\iota_{E'}(u)$ for the full logarithmic array with values
$(u_C,1-u_C)$ at each matched pair in a contracted component $C$.

We use exactly the fixed integer-spin Wilson ribbon representatives
declared in Paper 51, Section ``Actual Wilson insertions and their
joint angular limit'': passages through the designated positive cap
disks, fixed exterior tree corridors, and no new heavy-heavy crossing.
Their joint central symbol is $P_W(a,b)$, a polynomial even under
reversal of each spherical axis. Let $d\Pi$ be spherical probability.
For positive parameters $\rho=(\rho_e)_{e\in E'}$, set
\[
 A_j(a;\rho)=a_j+\sum_{(i\to j)\in E'}\rho_{ij}a_i.
\]
If $A\ne0$, let $d\sigma_A$ be normalized circle measure on
$\{b\in S^2:b\cdot A=0\}$. Define
\begin{equation}\label{eq:ratio-measure}
 d\nu_\rho(a,b)=\prod_j d\Pi(a_j)
       \prod_j\frac{d\sigma_{A_j(a;\rho)}(b_j)}{|A_j(a;\rho)|}.
\end{equation}
The zero sets of the $A_j$ have zero product-sphere measure; the
formula denotes its integrable extension, with no additional atoms.

\begin{lemma}\label{lem:ratio-finite}
The measure $\nu_\rho$ is finite and positive, and
$\nu_\rho(1)\le1$ for every positive $\rho$.
\end{lemma}
\begin{proof}
Choose a sink $j$ of the selected arrow forest. Then $a_j$ appears
only in $A_j$, with coefficient one. For a fixed vector $V$,
direct polar integration gives
\[
 \int_{S^2}\frac{d\Pi(a)}{|a+V|}
   =\frac1{\max(1,|V|)}\le1.
\]
Integrating $a_j$ and deleting its selected arrows bounds the
remaining integral by the same expression on a smaller forest.
Iteration proves the assertion. In particular all possible
$A_j=0$ singularities are integrable.
\end{proof}

\subsection{The quantum transition theorem}
At each vertex choose a fixed smooth profile $g_v(m_v/r)$ equal to
one near zero and supported in $[0,c)$, with $c<1/2$. These ordinary
negative-center profiles are retained. Let $F$ be a fixed smooth
joint test of all physical logarithmic radii $\log m_v/\log r$.
For $e=(i\to j)\in E'$, choose
$\chi_e\in C_c^\infty((0,\infty))$ and insert the literal
finite-level mark $\chi_e(m_{a_i}/m_{a_j})$. No product annulus or
Abel factor is part of this marked invariant.

\begin{theorem}[Physical ratio-resolved Wilson law]
\label{thm:ratio}
For the routes and marks just specified, put
$c_G=\sqrt2\,\pi^{1-N}$. Then
\begin{align}
 Z_{r,W}[F,g,\chi]
 &=(-1)^{r(N-1)}c_G r^{N-3/2}(\log r)^p
 \bigg[
 \int_{\Omega_{E'}}F(\iota_{E'}(u))\,d^pu\notag\\
 &\qquad\times
 \int_{(0,\infty)^{|E'|}}
      \left(\int P_W\,d\nu_\rho\right)
      \prod_{e\in E'}\chi_e(\rho_e)\frac{d\rho_e}{\rho_e}
 +o(1)\bigg].\label{eq:ratio-quantum}
\end{align}
The statement holds along the full integer sequence $r\to\infty$,
with $G,W$ and all profiles fixed. The profiles $\chi_e$ may be complex.
No uniformity as their ratio supports approach zero or infinity is
asserted.
\end{theorem}

\begin{proof}
Apply the finite-rank transfer expansion of Paper 51 and condition on
all $a$ labels. Sum the $b$ labels first, term by term. At $b_j$ let
$M_j$ be the largest neighboring $a$ label. The signed prefix bound
gives $Cr/M_j\le Cr/m_{a_j}$. Joint smooth logarithmic tests are
handled by their anchored mixed-variation representation in rectangular
prefixes, as in Paper 51. The ordinary profiles have uniformly bounded
variation on the retained lower branch.

The ratio marks confine all anchors in a selected component to fixed
positive multiples of one root. Successive leaf sums, or the change
to one root and one edge log ratio, give
\begin{equation}\label{eq:ratio-count}
 \sum_{\boldsymbol m_A\ {\rm odd}}
    \prod_jm_{a_j}^{-1}
    \prod_{(i\to j)\in E'}
       |\chi_{ij}(m_{a_i}/m_{a_j})|
       =O((1+\log r)^p).
\end{equation}
This lower logarithmic power is built into every subsequent bound.

Temporarily put $w_j=m_{a_j}m_{b_j}/r$. On $w_j>K\ge1$ the prefix
threshold is $K M_j/m_{a_j}$, so the conditional tail is bounded by
\[
 \frac{Crm_{a_j}}{KM_j^2}\le\frac{Cr}{Km_{a_j}}.
\]
Keep all other partner sums signed and use \eqref{eq:ratio-count}.
This costs $O(K^{-1})$ after normalization at the power in
\eqref{eq:ratio-quantum}. The estimate is uniform with arbitrary
additional hard interval restrictions on every other matched product:
at fixed anchors these are precisely partner-label intervals in the
prefix lemma. The ratio parameters are integrated against the fixed
$\chi$ in this assertion; a pointwise-in-$\rho$ finite-level estimate
is not needed. Similarly $w_j<\kappa$ contains
$O(\kappa r/m_{a_j})$ labels whenever nonempty. Bounded categorical
factors, with signed sums at the other partners, give $O(\kappa)$.
Thus all matched products may temporarily be kept in $[\kappa,K]$.

An unselected order violation $m_{a_i}/m_{a_j}>r^\epsilon$ gains
$r^{-\epsilon}$ from the $j$-th conditional bound. A width-$\epsilon$
strip about an unselected order wall has harmonic mass
$O(\epsilon(\log r)^p+(\log r)^{p-1})$. Restrict each selected
component root to $[r^\delta,r^{1-\delta}]$; the discarded mass has
the analogous bound with $\delta$. In the retained strict chambers,
all labels tend to infinity and all $m_v/r$ tend to zero. An
unselected product tends to zero, and a selected product is exactly
$\rho_{ij}w_j$. Paper 51's actual joint symbol theorem consequently
gives, uniformly on compact product boxes,
\[
 \mathcal K_{r,W}\longrightarrow
 \int P_W(a,b)\exp\!\left(i\pi\sum_j
               w_j b_j\cdot A_j(a;\rho)\right)
                  \prod_j d\Pi(a_j)d\Pi(b_j).
\]
The retained profiles $g_v$ are then one and $F$ tends to
$F(\iota_{E'}(u))$.

Each odd partner mesh gives $r\,dw_j/(2m_{a_j})$. The anchor harmonic
measure gives $2^{-d}(\log r)^p d^pu\prod_e d\rho_e/\rho_e$;
the log-incidence Jacobian has modulus one because $E'$ is a forest.
The $2d$ cap profiles tend to $2$ each. The exact common scalar
prefactor of the normalized sum in Paper 51 is asymptotic to
\[
 (-1)^{r(N-1)}2^{d+1/2}\pi^{1-2d}r^{d-3/2}.
\]
Thus the coefficient before the $w$ integration is $2^d c_G$.
This also follows from Paper 51's bounded-product overlap theorem,
by setting its selected product $v_{ij}=\rho_{ij}w_j$ and approximating
$\chi_{ij}(v_{ij}/w_j)$ on compact annuli by finite tensor sums.

Since $P_W$ is even in each $b_j$, the exponential can be replaced
by the product of cosines. Passing the interval-uniform bounds to
the compact-product limit gives $O(K^{-1})$ tails on arbitrary
rectangles, after integration against the fixed ratio profiles.
Multivariable summation by parts, or its continuous version, therefore
identifies the ordinary rectangular improper $w$ integral with its
Abel limit. For the Abel integral the following calculation is
pointwise for almost every axis configuration and ratio:
For $A\ne0$,
\[
 \int_0^\infty e^{-\varepsilon w}\cos(\pi ws)\,dw
      =\frac{\varepsilon}{\varepsilon^2+\pi^2s^2}
      \longrightarrow\delta_0(s),\qquad
 \delta(b\cdot A)d\Pi(b)=\frac{d\sigma_A(b)}{2|A|}.
\]
The spherical Abel mass is at most $(2|A|)^{-1}$. Lemma
\ref{lem:ratio-finite} therefore supplies an integrable dominating
function, also after integrating over the fixed compact ratio support.
The full $w$ integral is $2^{-d}\int P_Wd\nu_\rho$. Exhaust the
temporary product boxes and the root and order strips. Multiplication
by $2^dc_G$ proves \eqref{eq:ratio-quantum}.
\end{proof}

\subsection{The same extension in an intrinsic Hall sector}
Use the intrinsic decomposition $S_*,T_*,I$ of Paper 51, with
$a=|T_*|$, $b=|S_*|$, $|I|=2d$, $d\ge1$, and $N=a+b+2d$.
Here $I$ is a perfectly matched forest. Put
\[
 p_G=\frac{3N-2(a+d)-3}{2},\qquad
 c_G=2^{(1-N)/2+2a+d}\pi^{1-N}.
\]
Choose $E'$ among its unmatched dormant arrows and define
$p,\Omega_{E'},\iota_{E'}$ and $\nu_\rho$ on the dormant variables
as above. Let $\eta$ be a fixed smooth test with compact support
in the interior of all $S_*$ ordinary radii. Retain negative-center
profiles on $T_*\cup I$. The fixed Wilson words obey the
exterior-compatible routing of Paper 51: they may join dormant
components but pierce only dormant positive caps and have no retained
linking with a macroscopic ribbon after splitting the forced sublink.

\begin{corollary}[Hall ratio law with compact macro support]
\label{cor:ratio-hall}
For these tests and the literal dormant anchor-ratio gates,
\begin{align*}
 Z_{r,W}[\eta,F,g,\chi]
 &=(-1)^{r(N-1)}c_G r^{p_G}(\log r)^p
 \bigg[
 \left(\int\eta\,dM_*\right)
 \left(\int_{\Omega_{E'}}F(\iota_{E'}(u))\,d^pu\right)\\
 &\qquad\times
 \int\left(\int P_W\,d\nu_\rho\right)
       \prod_e\chi_e(\rho_e)\frac{d\rho_e}{\rho_e}+o(1)\bigg],
\end{align*}
where $dM_*$ is the positive active density already defined and
normalized in Paper 51. No matched-product or Abel gate is retained.
\end{corollary}
\begin{proof}
The pointwise forced estimate in Paper 51,
equation \texttt{eq:hall:uniform-forcing}, bounds the remaining
anchor summand after the signed partner sums by
\[
 C\prod_{t\in T_*}n_t^{-1-1/a}\prod_{j\in A_I}n_j^{-1}
 \qquad (a>0).
\]
Multiplication by the selected dormant ratio marks and
\eqref{eq:ratio-count} replace the total dormant logarithmic bound
by $(1+\log r)^p$. Consequently restricting every forced color to
$n_t\le H$ still has normalized tail $O(H^{-1/a})$, now at the
lower logarithmic order. The alternating paths used in this bound
lie entirely in the forced part, so they do not consume any dormant
partner estimate. For $a=0$ there is no forced exhaustion.

At fixed $H$, the same product-tail, root-strip, and order-strip
arguments as in Theorem~\ref{thm:ratio} apply. In a compact product
cell, the forced--dormant monodromy differs from the identity by
$O((Hm+H^2)/r)$, as proved in Paper 51. Its stated routing then splits
off the undecorated forced sublink while preserving the full joint
dormant symbol. The macro Riemann limit and subsequent $H\to\infty$
limit give $\int\eta\,dM_*$; compact macro support permits the
absolute character convergence already established there.

The finite compact-product normalization is $2^dc_G$. The dormant
Fourier integration contributes $2^{-d}\int P_Wd\nu_\rho$ exactly
as before. The uniform forced tail removes $H$ and proves the formula.
This corollary does not remove the compact macro test or the declared
negative-center profiles, and it does not add active cap passages.
\end{proof}

\subsection{An explicit nonzero next-logarithm Wilson signal on \texorpdfstring{$P_4$}{P4}}
Take the path $b_1-a_1-b_2-a_2$ and select its unmatched arrow
$a_1\to b_2$. Write $\rho=m_{a_1}/m_{a_2}$ and set
\[
 x=a_1\cdot b_1,\qquad y=a_1\cdot b_2,\qquad z=a_2\cdot b_2.
\]
Under the product sphere law these tree-edge dots are independent
and uniform on $[-1,1]$. Formula \eqref{eq:ratio-measure} is four
times that law restricted by $\delta(x)\delta(z+\rho y)$.
Its projection to $z$, with $a=\min(1,\rho)$, is consequently
\begin{equation}\label{eq:ratio-p4-law}
 x=0,\quad y=-z/\rho,\qquad
 d(\nu_\rho)_z=\frac1{2\rho}\mathbf1_{[-a,a]}(z)\,dz,
 \qquad \nu_\rho(1)=\frac1{\max(1,\rho)}.
\end{equation}
This is a full joint law with a coherent linear constraint.

For the actual spin-one word through caps $a_2,b_2$, the symbol is
$W_z=4z^2-1$. Hence
\[
 H(\rho):=\int(W_z+1)d\nu_\rho
 =\frac{4\min(1,\rho)^3}{3\rho}
 =\begin{cases}4\rho^2/3,&0<\rho\le1,\\
                 4/(3\rho),&\rho\ge1.
   \end{cases}
\]
With $F=1$ and any nonzero, nonnegative
$\chi\in C_c^\infty((0,\infty))$, Theorem~\ref{thm:ratio} gives
\begin{equation}\label{eq:ratio-nonzero}
 Z_{r,W_z}[g,\chi]+Z_r[g,\chi]
 =(-1)^r\frac{\sqrt2}{\pi^3}r^{5/2}\log r
   \left[\int_0^\infty\chi(\rho)H(\rho)\frac{d\rho}{\rho}
           +o(1)\right].
\end{equation}
The coefficient in brackets is strictly positive. Thus a literal
marked quantum observable has a nonzero next-logarithm-order angular
signal outside the principal perpendicular-pair support.

\begin{corollary}[A fixed hard ratio window]\label{cor:ratio-window}
Keep the declared negative-center profiles and restrict the actual
odd labels by $1\le m_{a_1}/m_{a_2}\le2$. With the same Wilson route
and framing in both terms, put
$D_r=Z_{r,W_z}[g,\mathbf1_{[1,2]}]+Z_r[g,\mathbf1_{[1,2]}]$.
Then
\[
 D_r=(-1)^r\frac{2\sqrt2}{3\pi^3}r^{5/2}\log r
                  +o(r^{5/2}\log r),
 \qquad
 Z_r[g,\mathbf1_{[1,2]}]
 =(-1)^r\frac{\sqrt2}{2\pi^3}r^{5/2}\log r
                  +o(r^{5/2}\log r).
\]
\end{corollary}
\begin{proof}
Approximate the window by smooth profiles differing only within
width $\eta$ of its endpoints. Conditional signed partner bounds
reduce their difference, for either Wilson term or the vacuum, to
the anchor harmonic mass of these strips. For $q=1,2$ this obeys
\[
 \sum_{\substack{m,n\ {\rm odd}\\|m/n-q|\le\eta}}
      \frac1{mn}
 \le C\sum_{n<r}\frac1n\left(\eta+\frac1n\right)
 =O(\eta\log r+1).
\]
The estimate also shows that an exact rational-ratio line has
harmonic mass $O(1)$ and therefore no atom at the logarithmically
normalized order. Let $r\to\infty$ and then $\eta\to0$.
Finally
$\int_1^2H(\rho)d\rho/\rho=2/3$ and
$\int_1^2\nu_\rho(1)d\rho/\rho=1/2$.
\end{proof}

The simultaneous spin-one tests $X=4z^2-1$ and $Y=4y^2-1$ further
give, after normalizing $\nu_\rho$, with $b=\min(1,1/\rho)$,
\[
 \mathbb EX=\frac{4a^2}{3}-1,\quad
 \mathbb EY=\frac{4b^2}{3}-1,\quad
 \operatorname{Cov}(X,Y)=\frac{64a^2b^2}{45}>0.
\]
In particular $\mathbb E z^2/\mathbb E y^2=\rho^2$; separate
independent angle marginals would lose this constraint.

\subsection{Every fixed integer spin}
The two-cap spin-$j$ character is $U_{2j}(z)$, with $U_n$ the
Chebyshev polynomial of the second kind. Its weighted coefficient is
\[
 J_j(\rho)=\frac1\rho\int_0^{\min(1,\rho)}U_{2j}(z)\,dz.
\]
For $\rho\ge1$ this is $1/[(2j+1)\rho]$; for $\rho\le1$ it is
an even polynomial in $\rho$. Subtract the principal perpendicular
value by defining
$H_j(\rho)=J_j(\rho)-(-1)^j/\max(1,\rho)$.
Then $H_j=O(\rho^2)$ at zero and $O(\rho^{-1})$ at infinity, and
\begin{equation}\label{eq:ratio-spin-integral}
 \int_0^\infty H_j(\rho)\frac{d\rho}{\rho}
 =\int_0^1\frac{U_{2j}(z)-(-1)^j}{z}\,dz
 =2(-1)^{j+1}\sum_{\substack{1\le k\le j\\k\ {\rm odd}}}\frac1k.
\end{equation}
Indeed absolute convergence permits exchanging the first two
integrals. If $I_j$ is the middle integral, the identity
$U_{2j}+U_{2j-2}=2zU_{2j-1}$ gives
$I_j+I_{j-1}=(1-(-1)^j)/j$, with $I_0=0$.

\paragraph{Remaining distinction.}
The integrability in \eqref{eq:ratio-spin-integral} alone does not
remove the ratio gates from the finite-level centered Wilson invariant.
That would require a bound preserving the cancellation between $W$
and the vacuum, uniform at the ratio endpoints and competing ordinary
radial endpoints. The separate prefix bounds of Paper 51 do not
establish that estimate. Nor has the present argument identified exact
color ratios with off-shell normal-form coordinates: it derives the
shifted angular constraints directly from the actual quantum sum.
